\section{Признак алгебраического замыкания}

\begin{proposition}[6.5]
Пусть $\Om/F$~--- расширение полей. Если $\Om$ алгебраично над $F$ и каждый непостоянный многочлен из $F[X]$ имеет
корень в $\Om$, то $\Om$ алгебраически замкнуто (и, значит, является алгебраическим замыканием $F$).
\end{proposition}

Условие слабее, чем в \rf{1.45}: там требовалось, чтобы каждый многочлен из $F[X]$ \emph{раскладывался} в $\Om$, а здесь~---
лишь чтобы у него был хотя бы \emph{один} корень. Про остальные корни ничего не известно; эту трудность обходят с
помощью примитивного элемента.

\begin{proof}
По \rf{1.45} достаточно доказать, что каждый унитарный неприводимый многочлен $f\in F[X]$ раскладывается в $\Om[X]$ на
линейные множители.

\emph{Шаг 1: случай сепарабельного $f$.} Пусть $E$~--- поле разложения $f$ над $F$. Расширение $E/F$ конечно и сепарабельно
(это расширение Галуа~--- поле разложения сепарабельного многочлена), поэтому $E=F[\gamma]$ по теореме о примитивном
элементе \rf{5.1}; именно здесь нужна сепарабельность. Пусть $g$~--- минимальный многочлен $\gamma$ над $F$. По условию у $g$
есть корень $\beta\in\Om$. Поля $F[\gamma]$ и $F[\beta]$~--- оба поля корня неприводимого многочлена $g$, поэтому есть
$F$-изоморфизм $\varphi\colon F[\gamma]\to F[\beta]\subset\Om$, $\gamma\mapsto\beta$ (\rf{2.1}). В $F[\gamma]$ многочлен $f$
раскладывается: $f=\prod_i(X-\alpha_i)$. Применяя $\varphi$ к коэффициентам (у $f$ они в $F$ и не меняются), получаем
$f=\prod_i\bigl(X-\varphi(\alpha_i)\bigr)$ в $\Om[X]$.
\[
\begin{tikzcd}[row sep=small,column sep=large]
 & E=F[\gamma]\arrow[dd,cnew,"\varphi"',"{\gamma\,\mapsto\,\beta}"] & \\
F\arrow[ur,hook]\arrow[dr,hook] & & \\
 & F[\beta]\arrow[r,hook] & \Om
\end{tikzcd}
\]
\figcap{В $F[\gamma]$ многочлен $f$ раскладывается на линейные множители; изоморфизм $\varphi$ (тождественный на $F$)
переносит это разложение в $F[\beta]\subset\Om$.}

В частности, предложение доказано, если $F$ совершенно: тогда все неприводимые многочлены над $F$ сепарабельны (\rf{2.24}).

\emph{Шаг 2: общий случай.} Осталось рассмотреть $\operatorname{char}F=p>0$. Положим
\[
F'=\{x\in\Om\mid x^{p^m}\in F\ \text{для некоторого}\ m\ge0\}.
\]
Докажем: (a) $F'$~--- подполе $\Om$, содержащее $F$; (b) $F'$ совершенно; (c) $\Om$ совершенно; (d) каждый непостоянный
многочлен из $F'[X]$ имеет корень в $\Om$. После этого шаг 1, применённый к полю $F'$ (над ним $\Om$ алгебраично, так
как алгебраично над $F$), показывает, что каждый многочлен из $F'[X]$ раскладывается в $\Om[X]$, и по \rf{1.45} поле $\Om$
алгебраически замкнуто.

(a) Ясно, что $F\subset F'$. Если $x^{p^m},y^{p^n}\in F$ и $k=\max(m,n)$, то $x^{p^k},y^{p^k}\in F$ (поле $F$ замкнуто
относительно возведения в степень $p$) и
\[
(x\pm y)^{p^k}=x^{p^k}\pm y^{p^k}\in F,\qquad (xy)^{p^k}\in F,\qquad (x^{-1})^{p^k}\in F\ \ (x\ne0),
\]
так как в характеристике $p$ возведение в степень $p^k$~--- гомоморфизм колец. Значит, $F'$~--- поле.

(b) Пусть $a\in F'$, $b=a^{p^m}\in F$. Многочлен $X^{p^{m+1}}-b\in F[X]$ имеет корень $\alpha\in\Om$, и
$\alpha^{p^{m+1}}=b\in F$, так что $\alpha\in F'$. При этом $(\alpha^p)^{p^m}=b=a^{p^m}$, а возведение в степень
$p^m$ инъективно ($x^{p^m}=y^{p^m}\Rightarrow(x-y)^{p^m}=0\Rightarrow x=y$). Значит, $\alpha^p=a$, и каждый элемент $F'$~---
$p$-я степень элемента $F'$.

(c) Пусть $a\in\Om$ не является $p$-й степенью в $\Om$. Тогда $X^p-a$ неприводим над $\Om$ (\rf{2.3}); пусть
$\Om[\alpha]$~--- его поле корня, $\alpha^p=a$. Рассмотрим цепочку $F'\subset F'[a]\subset F'[\alpha]$ внутри $\Om[\alpha]$.
Элемент $a$ не является $p$-й степенью и в $F'[a]\subset\Om$, поэтому $X^p-a$ неприводим над $F'[a]$ и
$[F'[\alpha]:F'[a]]=p$. Пусть $g$~--- минимальный многочлен $a$ над $F'$, $d=\deg g$. Тогда $[F'[\alpha]:F']=pd$, а
$g(X^p)$~--- унитарный многочлен степени $pd$ с корнем $\alpha$; значит, это минимальный многочлен $\alpha$ над $F'$, и
он неприводим. Но если $g=\sum c_iX^i$, то $c_i=e_i^p$ с $e_i\in F'$ по (b), и
$g(X^p)=\sum e_i^pX^{ip}=\bigl(\sum e_iX^i\bigr)^p$~--- противоречие с неприводимостью. Итак, $\Om$ совершенно.

\begin{center}
\begin{tikzpicture}
\node (A) at (0,0) {$F'$};
\node at (1.0,0) {$\subset$};
\node (B) at (2.1,0) {$F'[a]$};
\node at (3.2,0) {$\subset$};
\node (C) at (4.4,0) {$F'[\alpha]$};
\draw[decorate,decoration={brace,mirror,amplitude=5pt},cold] (-0.3,-0.3) -- node[below=6pt,font=\footnotesize]{$d=\deg g$} (2.6,-0.3);
\draw[decorate,decoration={brace,mirror,amplitude=5pt},cold] (2.7,-0.3) -- node[below=6pt,font=\footnotesize,align=center]{$p$ ($X^p-a$\\ неприводим)} (5.0,-0.3);
\draw[decorate,decoration={brace,amplitude=5pt},cnew] (-0.3,0.3) -- node[above=6pt,font=\footnotesize]{$pd$, поэтому $g(X^p)$~--- минимальный многочлен $\alpha$} (5.0,0.3);
\end{tikzpicture}
\end{center}
\figcap{Степени в доказательстве (c): если бы $a$ не был $p$-й степенью в $\Om$, то над совершенным полем $F'$
нашёлся бы неприводимый многочлен $g(X^p)$, являющийся $p$-й степенью.}

(d) Пусть $f=\sum_ia_iX^i\in F'[X]$ непостоянен. Коэффициентов конечное число, поэтому найдётся $m$ с $a_i^{p^m}\in F$
для всех $i$. Многочлен $\tilde f=\sum_ia_i^{p^m}X^i\in F[X]$ непостоянен и по условию имеет корень $\alpha\in\Om$. По
(c) $\alpha=\beta^{p^m}$ для некоторого $\beta\in\Om$, и
\[
f(\beta)^{p^m}=\sum_ia_i^{p^m}\beta^{ip^m}=\tilde f(\alpha)=0,
\]
т.\,е.\ $f(\beta)=0$.
\end{proof}
